% =====================================================================
%  Lecture 17 -- 3GPP Standardisation                     (English edition)
% =====================================================================
\documentclass[11pt,a4paper]{article}
\providecommand\ACNroot{../..}
\input{\ACNroot/templates/lecture-en}

\ACNset{number}{17}
\ACNset{title}{3GPP Standardisation}
\ACNset{shorttitle}{3GPP standardisation}
\ACNset{titleother}{استانداردسازی 3GPP}
\ACNset{part}{Part III -- Software-driven architectures (Lectures 15--18)}
\ACNset{time}{5--7 hours (3 h reading, 2--4 h exercises)}
\ACNset{level}{Introductory to intermediate}
\ACNset{prereq}{Lectures 1--16 (as examples); basic probability}
\ACNset{gen}{5}

\begin{document}
\ACNcover

% ---------------------------------------------------------------------
\section{Learning objectives}
After studying this lecture you should be able to:
\begin{objectives}
  \item \textbf{Describe} how 3GPP is organised: organisational partners, TSGs and working groups, and its
        relation to ITU-R and the IETF.
  \item \textbf{Explain} the three-stage methodology, specification numbering and the release model.
  \item \textbf{Decode} a specification version number and \textbf{classify} change requests.
  \item \textbf{Trace} how a feature moves from a study item to a specification and to products.
  \item \textbf{Locate and read} a specification efficiently: scope, normative text, annexes.
  \item \textbf{Reason quantitatively} about interoperability, work backlogs and time to market.
\end{objectives}

% ---------------------------------------------------------------------
\section{Prerequisite knowledge}
This lecture is about the process behind everything in this course. The earlier lectures provide the
examples: the specifications cited there (TS~23.401, TS~29.272, TS~38.300, \ldots) are the products of the
process described here. Exercise~5 uses Little's law (Lecture~15).

\paragraph{Key terms} The vocabulary of this lecture with its Persian equivalents:
\begin{keyterms}
3GPP & \fa{پروژهٔ مشارکت نسل سوم} & Partnership of standards organisations that specifies mobile systems. \\
Organisational partner & \fa{شریک سازمانی} & Regional/national standards body that is a member of 3GPP. \\
TSG / WG & \fa{گروه مشخصات فنی / گروه کاری} & Decision-making groups and their technical sub-groups. \\
TS / TR & \fa{مشخصات فنی / گزارش فنی} & Normative specification / informative study report. \\
Stage 1, 2, 3 & \fa{مراحل ۱، ۲، ۳} & Requirements, architecture, detailed protocols. \\
Release & \fa{انتشار} & Consistent set of features and specifications frozen together. \\
Freeze & \fa{انجماد} & Point after which only corrections are accepted. \\
Change request (CR) & \fa{درخواست تغییر} & Formal proposal to modify a specification under change control. \\
Work item / study item & \fa{مورد کاری / مطالعه} & Approved task to specify / to study a feature. \\
SEP / FRAND & \fa{پتنت ضروری استاندارد} & Standard-essential patent / fair, reasonable and non-discriminatory licensing. \\
\end{keyterms}

% ---------------------------------------------------------------------
\section{Introduction and motivation}
A phone bought in one country works on a network built from equipment of several vendors in another country.
That this works at all is the result of standards: thousands of documents that specify every message, every
timer and every bit, written by competitors who cooperate in the same rooms. For mobile networks the dominant
body is the \emph{3rd Generation Partnership Project} (3GPP), created in 1998 to specify the third generation and
since responsible for GSM's evolution, LTE, 5G and, now, 6G.

Engineers in the mobile industry spend much of their time with 3GPP specifications: implementing them, testing
against them, or contributing to them. Understanding how the documents are organised, how they evolve and how
decisions are made is therefore a practical skill, not trivia. It also explains why the systems look the way
they do: many design choices in this course are compromises reached in standardisation meetings.

% ---------------------------------------------------------------------
\section{Fundamental concepts}

\subsection{Organisation}
3GPP is not a legal entity of its own: it is a partnership of seven \emph{organisational partners} -- ARIB and TTC
(Japan), ATIS (USA), CCSA (China), ETSI (Europe), TSDSI (India) and TTA (Korea) -- which transpose its specifications
into their own standards. Companies participate as members of one of these partners; \emph{market representation
partners} (industry associations) contribute requirements. The technical work is done in three \emph{Technical
Specification Groups} and their working groups (\cref{fig:org}) \cite{tr21900}:
\begin{itemize}
  \item \textbf{TSG RAN}: radio access (physical layer, radio protocols, RAN architecture, RF, UE conformance tests);
  \item \textbf{TSG SA}: service and system aspects (requirements, architecture, security, codecs, management) and overall
        coordination;
  \item \textbf{TSG CT}: core network and terminals (NAS, core protocols, policy, USIM).
\end{itemize}
Working groups prepare documents; TSG plenaries, meeting quarterly, approve them. Decisions are made by consensus;
formal voting is possible but rare.

\begin{acnfigure}
  \resizebox{\linewidth}{!}{\acnfig{lecture-17/organisation}}
  \caption{Organisation of 3GPP: organisational partners, project co-ordination, TSGs and their working groups (selection).}
  \label{fig:org}
\end{acnfigure}

\subsection{Relation to other bodies}
ITU-R defines the requirements of each IMT generation (Lecture~14) \cite{itu2017m2410,itu2023m2160}; 3GPP's
technologies are submitted, through the organisational partners, as candidate IMT radio interfaces and evaluated.
The IETF specifies the Internet protocols that 3GPP reuses -- SCTP, Diameter, HTTP/2, SIP, IPsec, TLS
\cite{rfc6733,rfc9113,rfc9260,rfc3261} -- and 3GPP requests extensions when needed. Industry bodies such as the
GSMA define operational guidelines (e.g.\ roaming, VoLTE profiles), and certification organisations test devices
against 3GPP conformance specifications.

\subsection{Three stages}
Each feature is specified at three levels (\cref{fig:stages}):
\begin{itemize}
  \item \textbf{Stage~1} describes the service from the user's point of view: requirements, without solutions
        (e.g.\ TS~22.261 \cite{ts22261}).
  \item \textbf{Stage~2} describes the architecture: functional entities, reference points and information flows
        (e.g.\ TS~23.501 \cite{ts23501}).
  \item \textbf{Stage~3} specifies protocols bit by bit: messages, information elements, encodings, state machines
        (e.g.\ TS~24.501 \cite{ts24501}, TS~38.331 \cite{ts38331}).
\end{itemize}
Conformance-test specifications then define how devices are tested against stage~3.

\begin{acnfigure}
  \resizebox{\linewidth}{!}{\acnfig{lecture-17/stages}}
  \caption{Three-stage methodology with typical specifications.}
  \label{fig:stages}
\end{acnfigure}

\subsection{Specification numbering}
Every specification has a number \emph{ss.nnn} whose series \emph{ss} indicates the subject \cite{tr21900}.
\begin{table}[htbp]
\centering\small
\caption{Selected 3GPP specification series.}\label{tab:series}
\rowcolors{2}{acnmist}{white}
\begin{tabularx}{\linewidth}{@{}l Y l@{}}
\toprule
\rowcolor{white}Series & Subject & Example \\
\midrule
21 & Requirements, vocabulary, working methods & TR 21.905 \cite{tr21905} \\
22 & Service aspects (stage 1)                 & TS 22.261 \\
23 & Technical realisation (stage 2)           & TS 23.401, 23.501 \\
24 & Signalling protocols UE--core (stage 3)   & TS 24.301, 24.501 \\
26 & Codecs                                    & -- \\
28, 32 & Management and charging              & TS 28.530 \\
29 & Core network protocols (stage 3)          & TS 29.274, 29.500 \\
31 & SIM/USIM                                  & -- \\
33 & Security                                  & TS 33.401, 33.501 \\
35 & Security algorithms                       & TS 35.206 \\
36 & LTE (E-UTRA) radio                        & TS 36.300, 36.331 \\
37 & Multiple radio technologies               & TS 37.340 \\
38 & NR radio                                  & TS 38.300, 38.331 \\
\bottomrule
\end{tabularx}
\end{table}
A \emph{technical specification} (TS) is normative: products must comply with it. A \emph{technical report} (TR)
is informative: it records studies, evaluations and options, often in the 8xx and 9xx ranges (e.g.\ TR~38.913
\cite{tr38913}, TR~38.901 \cite{tr38901}).

% ---------------------------------------------------------------------
\section{Detailed technical explanation}

\subsection{Releases}
3GPP publishes its work in \emph{releases}: consistent sets of features specified together (\cref{fig:releases}).
Each release has freeze dates: after the \emph{stage~3 freeze} no new functions are added, only corrections; after the
\emph{ASN.1 freeze} the encoding of radio protocol messages is stable so that implementations become
interoperable. Releases overlap: while one is being corrected, the next is under development and the one after
is being studied. Vendors and operators implement a subset of features of a release, and devices declare their
capabilities, so ``supporting Release~17'' means little without a feature list.

\begin{acnfigure}
  \resizebox{\linewidth}{!}{\acnfig{lecture-17/releases}}
  \caption{Selected 3GPP releases with headline features. Years indicate approximate completion; teal marks 5G releases.}
  \label{fig:releases}
\end{acnfigure}

\subsection{Versions}
Each specification has a version \emph{x.y.z} \cite{tr21900}. Before approval, $x = 1$ (presented for information)
or $x = 2$ (presented for approval). Once approved and under change control, $x$ equals the release (e.g.\ 17 for
Release~17); $y$ increases with every set of approved technical changes; $z$ increases for purely editorial
changes. A specification thus exists in parallel versions, one per release, and corrections are often mirrored
into several releases.

\subsection{Change requests}
A specification under change control changes only through \emph{change requests} (CRs), each with a category
\cite{tr21900}: \textbf{F} (correction of an error), \textbf{A} (mirror of a category-F correction in a later
release), \textbf{B} (addition of a feature), \textbf{C} (functional modification of a feature) and \textbf{D}
(editorial). After freeze only F, A and D are normally accepted. CRs are agreed in the working group and approved by
the TSG; the specification is then republished with an incremented version (\cref{fig:cr}).

\begin{acnfigure}
  \resizebox{\linewidth}{!}{\acnfig{lecture-17/cr-process}}
  \caption{From study item to specification, and the change-request cycle.}
  \label{fig:cr}
\end{acnfigure}

\subsection{How work is done}
Companies submit written contributions (temporary documents, ``Tdocs'') to working-group meetings, held several
times per release cycle in different regions, with electronic meetings in between. Proposals are discussed,
merged and revised until consensus emerges; chairs record agreements. Typical contributions are proposals for
solutions in a study, text proposals for a TR or TS, and CRs. All documents are public on the 3GPP website, which
makes the history of every design decision traceable.

\subsection{Reading a specification}
A 3GPP specification follows a fixed template: scope, references, definitions and abbreviations, then the
technical clauses and annexes. Normative text uses ``shall'' (mandatory), ``should'' (recommended) and ``may''
(permitted); informative annexes give background and examples. Practical advice: read the stage-2 description
first, then the stage-3 procedure, and only then the message and information-element definitions; check the
release and version; and search across specifications, because a procedure is usually split over several of
them (e.g.\ an attach involves TS~23.401, 24.301, 36.413, 29.274 and 29.272).

\subsection{Intellectual property}
Companies that contribute technology often hold patents essential to implementing a standard (\emph{standard-essential
patents}, SEPs). Under the IPR policies of the organisational partners, they declare such patents and commit to
license them on fair, reasonable and non-discriminatory (FRAND) terms. Licensing is negotiated bilaterally, outside
3GPP, and the level of aggregate royalties is a recurring debate (\cref{sec:stack}).

% ---------------------------------------------------------------------
\section{Mathematical models}
The models below are deliberately simple; their numerical examples use illustrative assumptions, not measured
3GPP statistics.

\subsection{Interoperability across interfaces}\label{sec:interop}
\paragraph{Model} An end-to-end service crosses $k$ multi-vendor interfaces. If each is implemented correctly with
probability $p$, independently, the service works with probability
\begin{equation}\label{eq:interop}
  P_{\mathrm{e2e}} = p^{k}.
\end{equation}
\paragraph{Example} $p = 0.98$, $k = 20$: $P_{\mathrm{e2e}} = 0.67$. To reach 0.9 with $k = 20$ requires
$p \ge 0.9^{1/20} = 0.995$.
\paragraph{Interpretation} The more interfaces a system has, the more precise the specifications and the more
thorough the testing must be -- the reason for stage-3 detail, conformance tests and interoperability test events.
It is also an argument for fewer, simpler interfaces.

\subsection{Work in progress}\label{sec:little}
\paragraph{Model} If CRs on a specification group arrive at rate $\lambda$ (per month) and each spends on average $T$
months from first submission to approval, Little's law gives the number of open CRs:
\begin{equation}\label{eq:little}
  N = \lambda\,T .
\end{equation}
\paragraph{Example} With $\lambda = 400$ CRs/month and $T = 3$ months, about 1200 CRs are open at any time. Halving
$T$ (e.g.\ with electronic meetings between plenaries) halves the backlog for the same throughput.

\subsection{Time to market}\label{sec:ttm}
\paragraph{Model} A feature passes through study, normative work, implementation and testing; phases can partly
overlap. The time from study start to commercial availability is
\begin{equation}\label{eq:ttm}
  T_{\mathrm{market}} = \sum_{i} T_i - \sum_{j} \Delta_j ,
\end{equation}
with phase durations $T_i$ and overlaps $\Delta_j$.
\paragraph{Example} Study 12 months, stage~2 12, stage~3 12, implementation 18, testing and certification 6, with 6
months of overlap between implementation and late stage~3: $T_{\mathrm{market}} = 54$ months, i.e.\ about 4.5 years
-- consistent with the gap of several years between the start of a generation's studies and its first commercial
networks.

\subsection{Royalty stacking}\label{sec:stack}
\paragraph{Model} A device price $P$ carries royalties to $N$ SEP holders. If each independently asks a rate $r_i$, the
total is $R = \sum_i r_i$; under an aggregate cap $R_{\max}$ allocated by each holder's share $s_i$ of essential patents,
\begin{equation}\label{eq:stack}
  R = \sum_{i=1}^{N} r_i, \qquad r_i^{\mathrm{cap}} = R_{\max}\, s_i, \quad \sum_i s_i = 1 .
\end{equation}
\paragraph{Example} Twenty holders each asking 1\,\% stack to 20\,\% of the device price; with a 5\,\% cap, a holder with
10\,\% of the essential patents receives 0.5\,\%. The economics of SEPs are contested; the model only shows why the
question arises.

% ---------------------------------------------------------------------
\section{Network architecture and diagrams}
The organisation of 3GPP mirrors the architecture of the systems it specifies: RAN working groups own the radio
interface and RAN interfaces (Uu, X2/Xn, S1/NG, F1), SA2 owns the system architecture and reference points, CT
groups own the core protocols and the UE--core NAS. A new feature that crosses these boundaries -- network slicing,
for example -- needs coordinated work items in several groups, with SA2's stage~2 as the common reference.

% ---------------------------------------------------------------------
\section{Step-by-step procedures}

\begin{procedurebox}{Introducing a new feature}\label{proc:feature}
\begin{steps}
  \item Companies propose a study item (SI) with objectives; the TSG approves it.
  \item The working groups study use cases and candidate solutions and document them in a TR, with evaluations and a
        conclusion.
  \item A work item (WI) is approved with a scope, target release and responsible groups.
  \item Stage~2 and stage~3 specifications are written or extended (new TS or CRs to existing ones) until the
        stage~3 and ASN.1 freezes.
  \item After freeze, corrections continue through category-F CRs; conformance tests are specified; products
        implement and certify the feature.
\end{steps}
\end{procedurebox}

\begin{procedurebox}{Finding the right specification}\label{proc:find}
\begin{steps}
  \item Identify the topic and stage: requirement (22), architecture (23), protocol (24, 29, 36, 38), security (33).
  \item Start from the stage-2 overview (e.g.\ TS~23.501 or TS~38.300) and follow its references.
  \item Choose the release that matches the implementation and download the latest version of that release.
  \item Read scope and definitions, then the procedure clause, then the message definitions; check the CR history for
        recent changes.
\end{steps}
\end{procedurebox}

% ---------------------------------------------------------------------
\section{Worked examples}

\begin{examplebox}{Decoding a version}\label{ex:version}
\textbf{Problem.} What do ``TS~38.331 V17.5.0'' and ``TS~23.501 V2.0.0'' tell you?

\textbf{Solution.} TS~38.331 is in series 38 (NR radio): the NR RRC protocol. V17.5.0: Release~17, under change control,
fifth set of technical changes after its first Release-17 version, no editorial-only update since. TS~23.501 V2.0.0: a draft
presented to the TSG for approval (before it became V15.0.0 in Release~15).
\end{examplebox}

\begin{examplebox}{Interoperability budget}\label{ex:interop}
\textbf{Problem.} A VoLTE call involves about 20 interfaces between UE, eNB, MME, gateways, PCRF and IMS nodes from
different vendors. What per-interface correctness is needed for 99\,\% end-to-end success?

\textbf{Solution.} $p \ge 0.99^{1/20} = 0.9995$. Field experience with early VoLTE, which needed extensive interoperability
testing and an industry profile to narrow the options, illustrates the point.
\end{examplebox}

\begin{examplebox}{Classifying CRs}\label{ex:cr}
\textbf{Problem.} Classify: (a) a wrong timer value in Release~16, fixed in Release~16; (b) the same fix copied into
Release~17; (c) a new optional information element enabling a feature; (d) a typo.

\textbf{Solution.} (a) F; (b) A; (c) B (only before the release's freeze); (d) D.
\end{examplebox}

% ---------------------------------------------------------------------
\section{Technology comparison}
\begin{table}[htbp]
\centering\small
\caption{Standards bodies relevant to mobile networks.}\label{tab:bodies}
\rowcolors{2}{acnmist}{white}
\begin{tabularx}{\linewidth}{@{}l Y Y@{}}
\toprule
\rowcolor{white}Body & Role & Typical output \\
\midrule
ITU-R  & Spectrum regulation and IMT requirements      & Recommendations and reports (e.g.\ M.2083, M.2410, M.2160) \\
3GPP   & Specification of mobile systems               & TS and TR (series 21--38) \\
IETF   & Internet protocols                            & RFCs (e.g.\ SCTP, Diameter, HTTP/2) \\
ETSI   & European standards; NFV, MEC industry groups  & ENs, group specifications (e.g.\ NFV) \\
ONF, O-RAN Alliance & Open interfaces for SDN and RAN  & Specifications (e.g.\ OpenFlow, O-RAN interfaces) \\
GSMA   & Operator association                          & Permanent reference documents, profiles \\
\bottomrule
\end{tabularx}
\end{table}

\paragraph{Consensus versus speed} Open, consensus-based standardisation produces specifications that many companies
can implement and that last for decades, at the cost of time and complexity: options are added to satisfy many
parties. Industry alliances and open-source projects move faster but with less formal commitment; modern
networks combine all of them.

% ---------------------------------------------------------------------
\section{Key takeaways}
\begin{takeawaybox}
\begin{itemize}
  \item 3GPP = seven organisational partners; work in TSG RAN, SA and CT and their WGs; consensus decisions.
  \item Stage 1 (requirements), stage 2 (architecture), stage 3 (protocols); TS normative, TR informative.
  \item Series numbers identify the topic (23 architecture, 24/29 protocols, 33 security, 36 LTE, 38 NR).
  \item Releases freeze in stages; versions $x.y.z$ = release, technical changes, editorial changes.
  \item Change requests (F, A, B, C, D) are the only way to change a specification under change control.
  \item Interoperability decays as $p^k$; specifications and testing must be precise.
\end{itemize}
\end{takeawaybox}

% ---------------------------------------------------------------------
\section{Exercises}

\begin{exercisebox}{1}{Which series?}\label{exr:series}
In which series would you look for: (a) the NAS protocol of 5G; (b) the 5G security architecture; (c) the LTE MAC protocol;
(d) the service requirements of 5G; (e) the GTP-U protocol?
\end{exercisebox}

\begin{exercisebox}{1}{TS or TR?}\label{exr:tstr}
Is a channel model used to evaluate candidate technologies normative? Where is it published, and why does it matter for
comparing simulation results?
\end{exercisebox}

\begin{exercisebox}{2}{Versions}\label{exr:version}
A vendor's documentation states compliance with ``TS~24.501 V16.5.0''. The operator's network implements V16.9.0. What
kinds of differences can exist between the two, and how would you find them?
\end{exercisebox}

\begin{exercisebox}{2}{Interfaces and correctness}\label{exr:interop}
Using \cref{eq:interop}, compute $P_{\mathrm{e2e}}$ for $p = 0.995$ and $k = 20$ and 50. How many interfaces can a service
cross at $p = 0.999$ if $P_{\mathrm{e2e}} \ge 0.95$ is required?
\end{exercisebox}

\begin{exercisebox}{3}{Backlog}\label{exr:backlog}
A working group handles 250 CRs per month with an average time to approval of 4 months. A new release adds a burst so that
arrivals rise to 350 per month while approval capacity stays at 250 per month. Using \cref{eq:little} and a simple fluid
argument, estimate the backlog after 6 months of the burst. What can the group do?
\end{exercisebox}

\begin{exercisebox}{3}{Overlapping phases}\label{exr:ttm}
With the phase durations of \cref{sec:ttm}, how much earlier could commercial products appear if implementation started
when stage~2 is complete (overlapping stage~3 fully)? What risk does this create, and how do the freeze dates mitigate it?
\end{exercisebox}

\subsection*{Solutions}
\addcontentsline{toc}{subsection}{Solutions}

\begin{solutionbox}{\ref{exr:series}}
(a) 24 (TS~24.501); (b) 33 (TS~33.501); (c) 36 (TS~36.321); (d) 22 (TS~22.261); (e) 29 (TS~29.281).
\end{solutionbox}

\begin{solutionbox}{\ref{exr:tstr}}
No: channel models are informative; they are published in TRs (e.g.\ TR~38.901). They matter because simulation results from
different companies are comparable only if they use the same channel model and assumptions, which the TR fixes.
\end{solutionbox}

\begin{solutionbox}{\ref{exr:version}}
Both are Release~16; V16.9.0 includes four more sets of technical changes -- mostly corrections (category F), possibly
clarifications that change behaviour in corner cases. Compare the CR lists between the versions (the change history
annex and the CR database) and check whether any affect implemented procedures.
\end{solutionbox}

\begin{solutionbox}{\ref{exr:interop}}
$0.995^{20} = 0.905$; $0.995^{50} = 0.778$. With $p = 0.999$: $k \le \ln 0.95/\ln 0.999 = 51.3$, i.e.\ at most 51 interfaces.
\end{solutionbox}

\begin{solutionbox}{\ref{exr:backlog}}
Initially $N = 250\times4 = 1000$ open CRs. During the burst the backlog grows by $350 - 250 = 100$ per month: after 6 months
about 1600, and by Little's law the time to approval rises to about $1600/250 = 6.4$ months. Options: more meeting time
(additional electronic sessions), prioritisation of essential CRs, postponing optional features to the next release.
\end{solutionbox}

\begin{solutionbox}{\ref{exr:ttm}}
Implementation would overlap the whole 12 months of stage~3 instead of 6: products about 6 months earlier (48 instead of 54
months). Risk: implementing a moving target -- later stage-3 changes force rework and early products may not interoperate.
Freeze dates (especially the ASN.1 freeze) mark when the specification is stable enough to implement without such risk.
\end{solutionbox}

% ---------------------------------------------------------------------
\section{Further reading}
TR~21.900 \cite{tr21900} describes 3GPP working methods, version numbering and CR categories; TR~21.905 \cite{tr21905} is the
vocabulary. The ITU-R documents on IMT-2020 and IMT-2030 \cite{itu2015m2083,itu2017m2410,itu2023m2160} show how 3GPP's work
connects to global requirements. Dahlman et al.\ \cite{dahlman2020nr} describe the standardisation history of NR.

\printbibliography[heading=subbibliography,title={\ACNstrRefs}]
\end{document}
